% status: FULL
\chapter{Assembling a Decoder That Generates Tokens}\label{ch:decoder}
\begin{ChapterIntent}
All the pieces now exist. We will connect them into a small decoder whose
entire computation fits in a readable program. It will produce logits,
choose a token and continue from that choice. Then we will add a cache and
prove that the cache changes the work, not the intended prediction.
\end{ChapterIntent}

\section{Choose a model small enough to inspect}
Our teaching model has a vocabulary of five token IDs, model width four,
two query heads, one key/value head, two layers and an FFN width of six.
Each head has width two. The weights are deterministic values generated
from a sine function, not values learned from language.

That last sentence matters. The program is a complete numerical decoder
for the stated architecture, but its output is not meaningful language.
A random or hand-filled decoder can test shapes, causal dependencies,
position conventions and state handling. It cannot demonstrate the
language quality of a trained model.

We keep the vocabulary as IDs $0$ through $4$ so that tokenization is not
another moving part in this experiment. The prompt is $[0,1,2]$. A
trained application would obtain IDs using the tokenizer bound to its
checkpoint, then decode selected IDs with the same vocabulary. The text
boundary from Chapter~\ref{app:tokens} still applies.

The program uses pre-normalized residual blocks, full-width split-half
RoPE, grouped-query causal attention and a bias-free SwiGLU FFN.
Normalization gains are fixed to one. The output matrix is separate from
the embedding matrix. These choices define this teaching model; they are
not defaults that a loader may impose on an arbitrary checkpoint.

\section{Walk through the shapes}
Embedding lookup turns three IDs into three rows of width four:
$\mathbf X\in\mathbb R^{3\times4}$. Each layer reads this residual
matrix and emits another matrix of the same shape.

The query projection produces four coordinates per row, interpreted as
two heads of width two. Key and value projections each produce two
coordinates per row, interpreted as one head. Thus both query heads
share the same key/value head, but they can produce different scores.
RoPE rotates each Q and K head using the row's logical position:
zero, one or two. It does not rotate the value.

Each head returns two attention-output coordinates. Concatenating the
two heads restores four coordinates. The output projection maps those
four back to model width four. Add that update to the incoming residual.
Normalize the result for the FFN, expand along gate and up branches to
width six, multiply them after activation, project back to four and add
the second update.

After the second layer, normalize the final residual. The vocabulary
projection maps each width-four row to five logits. For next-token
generation we need the last row, the one associated with input token 2.
Earlier rows describe predictions at earlier prefix boundaries. They are
useful for tests, but they are not extra newly generated tokens.

\EngineFigure{foundation-decoder-cache}{Two executions of the same
illustrative prefix. Full-prefix evaluation rebuilds earlier intermediate
values. Cached evaluation keeps each layer's positioned K/V and processes
one newly supplied token. A selected token lies beyond the evaluated
frontier until feedback evaluates it.}{fig:foundation-decoder-cache}

\section{A complete prediction is not a complete generation}
Run the program from the repository root:
\begin{verbatim}
python3 code/foundations/lesson.py decoder
\end{verbatim}
The last logit row for $[0,1,2]$ is approximately
\[
 [0.093927,\ -0.204837,\ -0.450096,\ -0.577788,\ -0.554559].
\]
These are illustrative computations from the committed weights, rounded
for printing. Greedy selection chooses ID 0, whose score is largest.
It does not choose the coordinate with the largest absolute magnitude.
The logits can instead be passed to the sampling policy developed earlier.

Now append ID 0 to the token history and evaluate $[0,1,2,0]$. The
weights remain fixed. The input history and positions changed, so the
new final logit row can differ. Repeat this process until a token budget,
end token or another stopping condition ends generation.

The model function does not own the policy for when to stop. The
generation loop owns it. Our demonstration has no designated end token;
the tests use a fixed number of steps. Giving every toy vocabulary an
end token without defining how its parameters treat that token would
obscure the boundary we are trying to teach.

\section{First implement the expensive version}
The \texttt{Decoder.full} method embeds every input token, runs every
layer over every row, and returns a logit row for every position.
For each layer it constructs projections for the whole sequence, then
uses only the permitted prefix for each attention row.

This version has no persistent state. Calling it twice with the same
tokens and unchanged weights computes the same function. Its inefficiency
is a virtue while learning: there is no hidden cache to explain a wrong
answer. We can inspect all positions at every layer.

If we generate several tokens this way, each new call repeats work for
old positions. Causality tells us those earlier representations have not
changed. The earlier position cannot depend on the newly appended token.
We can exploit that fact, but we first need to identify exactly which
values a later query needs.

\section{Then keep only the reusable state}
A future query needs the earlier keys and values at each layer.
It does not need earlier query vectors. The \texttt{Cache} object
therefore owns one key history and one value history per layer, together
with the number of tokens already evaluated.

To append a new token, embed it and process one row through the stack.
At layer zero, normalize the row, calculate Q/K/V, rotate Q and K at
the current logical position, and append the new K/V to layer zero's
history. Attention reads that entire permitted history, including the
new entry. Apply the output projection, residual, normalized FFN and
second residual. Repeat at the next layer using its own cache.

The layers cannot share a K/V array merely because their shapes match.
Their weights differ, and each layer receives a different representation.
Likewise, two requests cannot share mutable cache histories merely
because their current lengths match. The state describes a particular
evaluated prefix under a particular model.

For the default model, each token adds two layers times two objects
(K and V) times one KV head times two coordinates: eight cached scalar
values. A packed F32 implementation would use 32 payload bytes per token.
Python lists and Python float objects take additional memory, so this
payload count is not a measurement of our program's resident size.

\section{Name the evaluated frontier}
After evaluating the three-token prompt, the cache position is three.
It contains representations for positions zero through two. Selecting
ID 0 does not change the cache position. Only evaluating that selected
ID advances the cache to four.

This distinction is easy to lose when a generation API returns a token
and a state object together. Ask whether the returned state includes
the returned token or only the prefix that predicted it. Either interface
can be designed consistently; an engine must not mix the two conventions.
Our interface is explicit: \texttt{append(token)} evaluates its argument,
then returns logits for the next selection.

RoPE uses this evaluated position, not the token ID, batch row or
number of calls to an unrelated request. Interleaving two requests must
not advance either one's position on behalf of the other. The tests
create two caches for the same model, interleave appends and compare
each result to its own full-prefix reference.

\section{Prove that reuse preserved the computation}
At every prefix length, compare the cached result with
\texttt{model.full(prefix)[-1]}. Compare all five logits, not only the
selected token. Two erroneous rows can still have the same maximum.
Use a stated tolerance; this example uses Python float arithmetic and
the tests allow an absolute difference of $10^{-12}$.

Run that comparison across different numbers of layers, different
histories, ordinary multi-head attention and grouped-query attention.
The committed tests do this for several small configurations. They
also change a future token and verify that earlier full-prefix rows
remain unchanged. Generation tests feed back the selected token for
several steps and compare both paths each time.

These are strong state and composition checks, but the paths share
scalar primitives. If both call a wrong softmax, their agreement will
not reveal it. That is why the previous chapters have hand-worked
attention and FFN fixtures, and why the test suite includes separate
normalization, linear-map and rotation checks.

Deliberately break the cached position by resetting it before each
append in a scratch copy. The first token still works, because its
correct position is zero. Later comparisons fail. A test containing
only singleton prompts would have missed the defect.

\section{What this decoder does not yet promise}
There is no checkpoint importer, training loop, optimized kernel,
scheduler or network server here. The cache is a growing Python list,
not paged device memory. It is bound to the decoder instance, whose
weights must remain unchanged during use. A general service would
need stronger lifecycle and mutation controls.

The existing Rust capstone implements a related decoder with paging,
batch scheduling and terminal events. It is useful later, after those
concepts have their own chapters. We do not make a beginner learn a
block allocator merely to understand a next-token prediction.
The small Python decoder is a transparent reference, not a new claim
that the early Rust ENGINE milestones already include these features.

\section{A model is now a concrete program}
You can now trace an ID into an embedding, through normalized residual
blocks, into logits, through selection and back into the next input.
You can also point to the state whose reuse makes incremental generation
possible. That is enough model knowledge to ask systems questions without
treating the Transformer as a black box.

The remaining chapters of Part I connect this computation to bytes,
numerical formats, hardware and experiments. They prepare the reader
to understand a roofline plot or a cache-capacity calculation as a
consequence of a program they already know, rather than a formula to
memorize.
\section{Worked problems}
\subsection*{1. The cache boundary}
\textbf{Problem.} A prompt has three tokens. The model has selected its
first continuation but has not evaluated that selection. How many positions
are cached?
\textbf{Solution.} Three, at positions zero through two. Evaluating the
selected token adds position three and advances the count to four.
\subsection*{2. Payload versus allocation}
\textbf{Problem.} Count the default teaching model's KV payload for five
evaluated tokens if stored as packed F32.
\textbf{Solution.} $2\cdot2\cdot1\cdot2\cdot4=32$ bytes per token,
so five tokens require 160 payload bytes. Python object overhead,
allocator metadata and any unused capacity are excluded.
\subsection*{3. A stronger oracle}
\textbf{Problem.} Why is agreement on the greedy token weaker than
agreement on every logit?
\textbf{Solution.} Many different score vectors have the same largest
entry. A faulty cache can change nonwinning scores while preserving that
winner. Compare complete rows at every prefix and include cases with
different token histories and head sharing.
\Needspace{8\baselineskip}
\subsection*{4. Count the actual model}
\textbf{Problem.} Count matrix parameters and hypothetical packed F32
payload for the default decoder, including the untied vocabulary output.
\textbf{Solution.} Each layer has $16+8+8+16+24+24+24=120$ entries.
Embedding and output each add twenty. Total is $20+240+20=280$,
or 1,120 payload bytes at four bytes per entry. Unit normalization gains
and absent biases add no learned parameters.
\Needspace{8\baselineskip}
\subsection*{5. Three predictions, two feedback evaluations}
\textbf{Problem.} Compare projected rows and causal query--key pairs when
predicting three IDs from a three-token prompt, across both layers and
both query heads of the default model.
\textbf{Solution.} Full-prefix lengths are three, four and five: projected
rows total $2(3+4+5)=24$, and pairs total $4(6+10+15)=124$.
Cached evaluation visits five tokens once: ten projected rows and
$4(1+2+3+4+5)=60$ pairs. These are work counts, not timings.
\Needspace{8\baselineskip}
\subsection*{6. Stop without evaluating the end ID}
\textbf{Problem.} Suppose the first selected ID is configured as the end
ID in the trace's policy. How many IDs are selected and how many positions
are cached?
\textbf{Solution.} One ID is selected and retained in the trace. The
three prompt positions remain cached; the selected end ID is not fed back.
Stopping and evaluation are different events.

